Los resultados obtenidos nos permiten evidenciar de forma experimental las diferencias prácticas entre los distintos algoritmos analizados. En el caso del ordenamiento, algoritmos como MergeSort y QuickSort mostraron un rendimiento significativamente superior a SelectionSort, tanto en tiempo como en escalabilidad, mientras que \texttt{std::sort}, optimizado a nivel de biblioteca, superó incluso a los algoritmos recursivos clásicos en eficiencia.

El análisis de uso de memoria, por otro lado, mostró que la eficiencia temporal no siempre va de la mano con un bajo uso de recursos. QuickSort, por ejemplo, mostró un pico inesperado en uso de memoria para entradas grandes, lo que resalta la importancia de considerar múltiples métricas al evaluar algoritmos.

Respecto a la multiplicación de matrices, los experimentos mostraron que el algoritmo de Strassen, si bien conceptualmente más eficiente para grandes entradas, incurre en un costo de memoria elevado. El método Naive, en comparación, resultó competitivo e incluso superior en tiempo, mostrando que su simplicidad lo vuelve una alternativa razonable en contextos con pocos datos.

En conjunto, los resultados obtenidos permiten reforzar el entendimiento de que la elección del algoritmo más adecuado depende no sólo de su complejidad teórica, sino también del tipo de entrada, del entorno de ejecución y de los recursos disponibles. Este análisis experimental complementa al estudio teórico y nos entrega una visión más completa sobre el rendimiento de los algoritmos en distintos escenarios.

